
\section{Method}
\label{sec:method}

 Over the course of its lifetime a capable \acrfull{rl} agent will exhibit complex behaviours, such as exploration, temporal credit assignment, and planning. Our key insight is that an agent's actions, regardless of the environment it inhabits, its internal structure, and implementation, can be viewed as a function of its past experience, which we refer to as its \emph{history}. Formally, we write:
\looseness=-1
%
\begin{align}
  \mathcal{H} \ni h_{t} \defeq \left( o_{0}, a_{0}, r_{0}, \ldots, o_{t-1}, a_{t-1}, r_{t-1}, o_{t}, a_{t}, r_{t} \right) = \left( o_{\leq t}, r_{\leq t}, a_{\leq t}  \right)
\end{align}
%
and we refer to a \emph{long}\footnote{%
  Long enough to span learning updates, \eg~across episodes.
  }
history-conditioned policy as an \emph{algorithm}:
\begin{align}
  \algo: \mathcal{H}  \cup \mathcal{O} \rightarrow \Delta(\mathcal{A}),
\label{eq:algo-signature}
\end{align}
%
where $\Delta(\mathcal{A})$ denotes the space of probability distributions over the action space $\mathcal{A}$.
\eqnref{eq:algo-signature}~suggests that, similar to a policy, an algorithm can be unrolled in an environment to generate sequences of observations, rewards, and actions.
For brevity, we denote the algorithm as $\algo$ and environment (\ie~task) as $\mathcal{M}$, such that the history of learning for any given task $\mathcal{M}$ is generated by the algorithm $\algo_{\mathcal{M}}$.
%
\begin{align}
  \left( O_{0}, A_{0}, R_{0}, \ldots, O_{T}, A_{T}, R_{T} \right) \sim \algo_{\mathcal{M}}.
\end{align}
Here, we're denoting random variables with uppercase Latin letters, \eg~$O$, $A$, $R$, and their values with lowercase Latin letters, \eg~$o$, $a$, $r$.
By viewing algorithms as long history-conditioned policies, we hypothesize that any algorithm that generated a set of learning histories can be distilled into a neural network by performing behavioral cloning over actions.
Next, we present a method that, provided agents' lifetimes, learns a sequence model with behavioral cloning to map long histories to distributions over actions.

\subsection{Algorithm Distillation}
\label{subsec:algorithm-distillation}

Suppose the agents' lifetimes, which we also call {\it learning histories}, are generated by the source algorithm $\algo^{\text{source}}$ for many individual tasks $\{ \mathcal{M}_{n} \}_{n=1}^{N}$, producing the dataset $\mathcal{D}$:
\looseness=-1
%
\begin{align}
  \mathcal{D} \defeq \left\{ \left( o_{0}^{(n)}, a_{0}^{(n)}, r_{0}^{(n)}, \ldots, o_{T}^{(n)}, a_{T}^{(n)}, r_{T}^{(n)} \right) \sim \algo_{\mathcal{M}_{n}}^{\text{source}} \right\}_{n=1}^{N}.
  \label{eq:data}
\end{align}
%
Then we distill the source algorithm's behaviour into a sequence model that maps long histories to probabilities over actions with a negative log likelihood (NLL) loss and refer to this process as \emph{algorithm distillation} (AD).
In this work, we consider neural network models $\algo_{\theta}$ with parameters $\theta$ which we train by minimizing the following loss function:
\looseness=-1
%
\begin{align}
   \mathcal L(\theta) \defeq  -\sum_{n=1}^{N} \sum_{t=1}^{T-1} \log \algo_{\theta}(A=a_{t}^{(n)} \vert h_{t-1}^{(n)}, o_t^{(n)}).
\label{eq:thetastar}
\end{align}
%
Intuitively, a sequence model with \emph{fixed} parameters that is trained with AD should amortise the source RL algorithm $\algo^{\text{source}}$ and by doing so exhibit similarly complex behaviours, such as exploration and temporal credit assignment. Since the RL policy improves throughout the learning history of the source algorithm, accurate action prediction requires the sequence model to not only infer the current policy from the preceding context but also infer the improved policy, therefore distilling the policy improvement operator.

\begin{algorithm}[h]
\caption{\texttt{Algorithm Distillation}} \label{alg:ad_full}
\begin{algorithmic}[1]
\footnotesize

\Require Train $\{ \mathcal M^{\text{train}} \}$ and test $\{ \mathcal M^{\text{test}} \}$ tasks, observations $o \in \mathcal O$, actions $a \in \mathcal A$, and rewards $r \in \mathcal R$.
\Require Network parameters $\phi_i$ for $i=1,\dots, N$ source RL algorithms. 
\Require Network parameters $\theta$ for a causal transformer $\algo_{\theta}$ that predicts actions.
\Require An empty buffer to store data $\mathcal D$.
\For{$i = 1\dots N$} \Comment{Part 1: Dataset Generation}
\State Sample a task $\mathcal{M}^{\text{train}}_i$ randomly from the train task distribution.
\State Train the source RL algorithm $\phi_i$ until it converges to the optimal policy.
%\If{successful}:
\State Save the learning history $h^{(i)}_{T} = \left( o_{0}, a_{0}, r_{0}, \ldots, o_{T}, a_{T}, r_{T} \right)_i$ to the dataset $\mathcal D \leftarrow \mathcal D \cup h^{(i)}_T $.
%\EndIf
\EndFor
\
\While{$P_\theta$ not converged} \Comment{Part 2: Algorithm Distillation}
\State Randomly sample a multi-episodic subsequence $\bar h^{(i)}_j = \left( o_{j}, a_{j}, r_{j}, \ldots, o_{j+c}, a_{j+c}, r_{j+c} \right)_i$ of length $c$.
\State  Autoregressively predict the actions with $\algo_\theta$ and compute the NLL loss in Eq.~\ref{eq:thetastar}.
\State Backpropagate to update the transformer parameters.
\EndWhile
\For{$k = 1\dots M_{\text{seeds}}$}  \Comment{Part 3: Autoregressive Evaluation}
\State Sample a task $\mathcal{M}^{\text{test}}_k$ randomly from the test task distribution. Initialize empty context queue $C$.
\State Unroll the transformer $P_\theta (\cdot | C)$ in the environment  storing sequential transitions (\ie\, histories) in $C$.
\State Measure the return accumulated by the agent for each episode of evaluation.
\EndFor
\end{algorithmic}
\end{algorithm}

\subsection{Practical Implementation}
\label{subsec:practical-implementation}

In practice, we implement \name as a two-step procedure. First, a dataset of learning histories is collected by running an individual gradient-based RL algorithm on many different tasks. Next, a sequence model with multi-episodic context is trained to predict actions from the histories. We describe these two steps below and detail the full practical implementation in Algorithm~\ref{alg:ad_full}.

\noindent {\bf Dataset Generation:} A dataset of learning histories is collected by training $N$ individual single-task gradient-based RL algorithms. To prevent overfitting to any specific task during sequence model training, a task $\mathcal{M}$ is sampled randomly from a task distribution for each RL run. The data generation step is RL algorithm agnostic - any RL algorithm can be distilled. We show results distilling UCB exploration~\citep{lai86ucb}, an on-policy actor-critic~\citep{mnih16a3c}, and an off-policy DQN~\citep{mnih2013dqn}, in both distributed and single-stream settings. We denote the dataset of learning histories as $\mathcal D$ in Eq.~\ref{eq:data}.
%
\noindent {\bf Training the Sequence Model:} Once a dataset of learning histories $\mathcal D $ is collected, a sequential prediction model is trained to predict actions given the preceding histories. We utilize the GPT~\citep{radford2018improving} causal transformer model for sequential action prediction, but \name is compatible with any sequence model including RNNs~\citep{williams1989learning}. For instance, we show in Appendix~\ref{app:transformer_vs_lstm} that \name can also be achieved with an LSTM~\citep{schm97lstm}, though less effectively than \name with causal transformers. Since causal transformer training and inference are quadratic in the sequence length, we sample across-episodic subsequences $\bar h_j = \left( o_{j}, r_{j}, a_{j} \ldots, o_{j+c}, r_{j+c}, a_{j+c} \right)$ of length $c<T$ from $\mathcal D$ rather than training full histories. 
